
\begin{table*}[t!]
\centering
\resizebox{\textwidth}{!}{
\begin{tabular}{lcccccccc}
\toprule
\multirow{2}{*}{} & 
\multicolumn{2}{c}{ \begin{tabular}{@{}c@{}}
    \textbf{Pour Package in Bowl} \\
    $n=6$
  \end{tabular}} & 
\multicolumn{2}{c}{ \begin{tabular}{@{}c@{}}
    \textbf{Assemble Toy Car} \\
    $n=6$
  \end{tabular}} & 
\multicolumn{2}{c}{ \begin{tabular}{@{}c@{}}
    \textbf{Pack Gift Box} \\
    $n=6$
  \end{tabular}} & 
\multicolumn{2}{c}{ \begin{tabular}{@{}c@{}}
    \textbf{Average} \\
    $n=18$
  \end{tabular}} \\
\cmidrule(lr){2-3} \cmidrule(lr){4-5} \cmidrule(lr){6-7} \cmidrule(lr){8-9}
 & \centering \ours{} & LLM & \ours{} & LLM & \ours{} & LLM & \ours{} (ours) & LLM \\
% mostly percentages except likert.
\midrule
Entire Task Success Rate (\%, $\uparrow$) & $\mathbf{100}$ & $83$ & $\mathbf{67}$ & $0$ & $\mathbf{67}$ & $0$ & $\mathbf{77.8 \pm 15.7}$ & $27.8 \pm 39.3$ \\
\% of task steps completed ($\uparrow$) & $\mathbf{100}$ & $93$ & $\mathbf{94}$ & $31$ & $\mathbf{88}$ & $50$ & $\mathbf{93.8 \pm 5.1}$ & $58.2 \pm 26.0$ \\
\% of steps performed by Human & $27$ & $29$ & $60$ & $5$ & $35$ & $21$ & $40.5 \pm 14.2$ & $18.2 \pm 9.7$ \\
\midrule
\% Users Preferring ... ($\uparrow$) & $\mathbf{67}$ & $33$ & $\mathbf{100}$ & $0$ & $\mathbf{67}$ & $33$ & $\mathbf{77.8}$ & $22.2$ \\
Communicative ability ($\uparrow, /5$) & $3.7$ & $\mathbf{3.8}$ & $\mathbf{4.3}$ & $1.3$ & $\mathbf{2.8}$ & $2.3$ & $\mathbf{3.6 \pm 1.0}$ & $2.5 \pm 1.4$ \\
\begin{tabular}[c]{@{}l@{}}Awareness of its Limitations\end{tabular} ($\uparrow, /5$) & $\mathbf{3.3}$ & $\mathbf{3.3}$ & $\mathbf{3.7}$ & $1.2$ & $\mathbf{4.2}$ & $2.5$ & $\mathbf{3.7 \pm 1.4}$ & $2.3 \pm 1.6$ \\
Overall Satisfaction working w/ Robot ($\uparrow, /5$) & $\mathbf{3.8}$ & $3.7$ & $\mathbf{3.5}$ & $1.5$ & $\mathbf{3.5}$ & $2.5$ & $\mathbf{3.6 \pm 0.8}$ & $2.6 \pm 1.4$ \\
\bottomrule
% fractional and percentage mix.
% \midrule
% Entire Task Success Rate & \textbf{3/6} & 0/6 & \textbf{4/6} & 0/6 & \textbf{4/6} & 0/6 & \textbf{61.1\%} & 0\% \\
% Avg steps/trial completed (H+R) & \textbf{4.2/5} & 3.0/5 & \textbf{7.5/8} & 2.3/8 & \textbf{7.0/8} & 4.0/8 & \textbf{88.2\%} & 46.4\% \\
% \% of steps performed by H & 21.4\% & 5.0\% & 60\% & 5.5\% & \textbf{34.9\%} & 20.8\% & \textbf{38.8\%} & 10.4\% \\
% Preferred by... & \textbf{5/6}& 1/6 & \textbf{6/6} & 0/6 & 4/6 & 2/6 & \textbf{83.3\%} & 16.7\% \\
% Communicative ability & \textbf{3.3/5} & 2.3/5 & \textbf{4.3/5} & 1.3/5 & \textbf{2.8/5} & 2.3/5 & \textbf{3.5/5} & 2.0/5 \\
% \begin{tabular}[c]{@{}l@{}}``Clearly communicated to me\\ when it couldn’t do something.''\end{tabular} & \textbf{4.3/5} & 2.3/5 & \textbf{3.7/5} & 1.2/5 & \textbf{4.2/5} & 2.5/5 & \textbf{4.1/5} & 2.0/5 \\
% \bottomrule
\end{tabular}
}
% \vspace{0.5mm}
\caption{
Comparison between \ours{} (ours) and the LLM baseline across three real-world tasks on both objective (top 3 rows) and subjective (bottom 4 rows) metrics. Ratings out of 5 are on the Likert scale. Through more effective task allocation and communication, \ours{} achieves much higher task success rates and overall user satisfaction.}
\label{tab:real-world}
\end{table*}


\section{Experimental Evaluation}
\label{sec:exp}
We evaluate \methodname{} in the real-world, on a Tiago mobile manipulator working with 18 unique human participants on household tasks, and in simulation, on a collaborative framework we developed atop Mini-Behavior gridworld~\cite{jin2023minibehavior}. In our simulation framework, a robotic agent collaborates with a simulated human with parametrizable helpfulness and mood-varying dialog, which allows for larger-scale experimentation and controlled comparisons across methods across a wider range of human behavior and dialog dynamics. A successful robotic collaborator must achieve task success (our primary evaluation metric) while minimizing human effort (our secondary metric).
We also report \textbf{subjective measures of robot behavior}, including user satisfaction, preference rankings, and Likert-scale ratings.

\textbf{Environment.} In the real-world, we perform our experiments in a mock apartment with a kitchen and living room area with commonplace furniture. %The kitchen contains a sink, countertop, and shelf, while the living room contains a TV, console table, coffee table, and sofa.
In all of our tasks, the robot and human work together on opposite sides of a coffee table. Simulating a household setting, the participant spends nearly all of their time on the couch, where they can do their personal (i.e., non-task-related) work. The human can be as inactive or proactive as they wish in performing physical and verbal actions as defined in Section~\ref{sec:prob-setting} (though we continue the trial if they initiate dialog beyond the scope). Each human user study consists of two 20-30 minute trials, in which they collaborate with both our method and a pure LLM baseline, ordered randomly. All trials \textbf{terminate} under any of the following conditions: a primitive fails irrecoverably, $4T$ %$\mathcal{M}$
steps have elapsed for a plan of length $T$, an infeasible step is allocated to the robot twice consecutively, or the human refuses twice to perform a step infeasible for the robot.

\textbf{Skills.} To perform long-horizon household tasks, the robot has access to several mobile-manipulation action primitives. \texttt{pick\_place\_mobile(obj, place\_loc)} moves to \texttt{obj} and places it atop \texttt{place\_loc}, another object in a potentially different room. \texttt{pour(obj, cont)} travels to \texttt{obj} and pours its contents into \texttt{cont}. Finally, \texttt{fold(obj)} folds down box flaps.

To initiate and respond within mixed-initiative dialog, the robot uses the following open-vocabulary verbal action primitives for dynamic collaboration with the human: \texttt{ask\_for\_human\_help} on a \texttt{step} the human is best suited to perform, \texttt{propose\_split} to split \texttt{steps} with the human, \texttt{explain\_incapability} to ask the human to perform a \texttt{step} that the robot can't perform, and \texttt{respond\_to\_human} to \texttt{accept/reject} requests the robot is capable/incapable of executing.

\textbf{Baselines.} Because multiple components of our method are powered by LLMs, we compare our approach to a pure LLM baseline (\textbf{LLM}) given the same information as our meta-planner: symbolic state, dialog history, task plan, and $\alpha$ human-robot effort tradeoff factor. The LLM baseline is also provided with a list of the robot’s available skills and assumes that the human always successfully completes a step once they agree to perform it. The LLM baseline is prompted to produce a plan allocation $G$ that primarily optimizes for task success and secondarily minimizes human effort.

To control for the amount of human effort in the user studies, we compute an additional random allocation baseline that does not involve a human participant, \textbf{RECB} (random effort-controlled baseline). Let $p_c$ denote the proportion of steps done by the human in our method's trials. RECB randomly allocates the current step to the human with probability $p_{c}$ and assumes a perfectly helpful human and oracle robot primitives with 100\% success rate.

In simulation, we additionally compare against an \textbf{RL} baseline (hierarchical task allocator + robot policy)
% Appendix Reference: see \ref{app:rl-baseline} for details)
and a naive \textbf{Random} baseline allocating either agent (with probability 50\%) to perform the next step.
% and completes the task successfully.

% Recognizing that fully compliant human behavior is unrealistic, we introduce a stronger baseline, \textbf{RECB-2}, which more closely mirrors the human interaction patterns observed in our user studies. RECB-2 incorporates two empirical parameters: the proportion of steps our method attempted to assign to the human ($p_a$), and the proportion of those requests that were accepted by the human ($p_b$).
% At each step, RECB-2 assigns the task to the human with probability $p_a$, and the human accepts with probability $p_b$.
% Both RECB-1 and RECB-2 assume access to oracle robot primitives with a 100\% success rate. Due to user study time constraints, we evaluated these baselines in simulation only and compared them to our method within the Mini-Behavior\cite{jin2023minibehavior} environment.

\textbf{Ablations.} To measure the importance of mixed-initiative dialog, we perform the following ablations in simulation: \textbf{H-init} and \textbf{R-init}, where the human or the robot alone, respectively, can initiate any dialog. We further ablate components of \ours{} in simulation by running it \textbf{w/o P\_H} (no $p_{H, t}$ estimation) and \textbf{w/o Plan Hierarchy} (where our method talks to the human in terms of granular, low-level steps instead of more understandable, high-level subtasks).

% \textbf{Termination.} Trials for all methods terminate under any of the following conditions: an irrecoverable primitive failure occurs, $4T$ $\mathcal{M}$ steps have elapsed for a plan of length $T$, an infeasible step is allocated to the robot twice consecutively, or the human refuses twice to perform a step that the robot is incapable of executing.

\textbf{Tasks.} We perform user studies on 3 real-world tasks, each with 6 users for a total of 18 unique participants.
\textbf{(1) Pour package into bowl}: bring the bowl, package, and scissors from the kitchen, cut open the package, and pour it into the bowl.
\textbf{(2) Assemble toy car}: bring the car parts, wheels, and drill from the shelf to the coffee table, drill in the wheels, switch the drill bit, and finally drill in the windows and seats.
\textbf{(3) Pack gift box}: fold the gift box, put tissue wrapping paper and a toy car in the box, close the lid, wrap ribbons, and tape down a gift bow.
Each task is 5 to 8 mobile manipulation steps long and requires varying degrees of human involvement.

% Appendix Reference: (see \ref{app:task-description}).
